% Round 20, inertia20. New independent verification of a historical R9
% candidate. Not a modification or automatic activation of the old archive.
\section{Two bad surface factors: a staircase and projector verification}
\label{r20px:section}

\paragraph{Status and provenance.}
The argument below independently reconstructs the two-factor part of
the historical R9 ST/TR candidate. The old files are not treated as
frozen theorems. The new surface-incidence inverse of AS20/PI20 is used
only after its elementary proof is recalled below. This note proves
the geometric component and incidence assertions needed for two bad
surface factors, including genuine triple intersections in either
factor. It does not activate the historical arbitrary-factor DT/UP
claim. Independent review of this new source is required before it
enters the Round 20 frozen theorem register.

All Chow groups have rational coefficients. A cohomological identity
of correspondences means an identity on full geometric etale
cohomology, with the actual degree and Tate shifts; it need not be an
identity in the Chow group. Numerical degrees below are always on
proper varieties over a finite residue field.

\subsection{The ordered staircase and its component towers}

Let $R$ be a complete DVR with finite residue field $k$, and let
$\mathcal S,\mathcal T$ be regular projective reduced strictly
semistable surface models. After a finite unramified extension, split
the finitely many component and stratum constants and the actions on
the geometric Neron--Severi groups of all surface components.
The zero-dimensional strata and all point centers in the surface
factorizations below may be split by a further finite unramified
extension; such extensions preserve the entire construction.
Fix total orders
on the components $S_i$ and $T_j$. Put
$E_{ia}=S_i\cap S_a$, $F_{jb}=T_j\cap T_b$, treating nonempty
connected components separately.

At a product stratum with $r$ and $s$ local components, the vertical
cone has lattice and rays
\[
 N=\{(a,b)\in\mathbf Z^r\oplus\mathbf Z^s:
              \textstyle\sum a=\sum b\},\qquad
 \sigma=N_{\mathbf R}\cap\mathbf R_{\geq0}^{r+s},\qquad
 n_{ij}=(e_i,f_j).
 \tag{PX20.1}
\]
Its height-one section is $\Delta^{r-1}\times\Delta^{s-1}$.
Subdivide it into the monotone-path simplices for the fixed orders.
The resulting cone complex glues on every face. Each maximal cone is
unimodular: one first ray, followed by successive path-ray differences,
is an integral basis of $N$. Every vertical ray still has height one.
Thus the resulting toroidal modification $\mathcal P$ is regular
strictly semistable and has precisely the components $P_{ij}$
continuing $S_i\times T_j$.

We specify projectivity, rather than silently impose it on an arbitrary
subdivision. Give the vertex $(i,j)$ height $-ij$, using the globally
chosen integer positions in the two orders. These heights give the
staircase regular subdivision. Indeed, if a positive convex
combination uses two incomparable pairs $(a,b'),(a',b)$ with
$a<a'$ and $b<b'$, transferring equal small mass to
$(a,b),(a',b')$ preserves both marginal coordinates and strictly
decreases height by a positive multiple of
$(a'-a)(b'-b)$. A minimizing combination therefore has chain support.
The monotone coupling of the two marginal distributions gives such
a chain, and its vertices extend to a monotone path. Conversely,
strictly positive masses on a maximal path, with distinct cumulative
marginals, give its unique monotone coupling; hence every path is a
lower simplex. This is the claimed lower subdivision. The integral
piecewise linear functions
agree on faces, so the toroidal modification is projective; equivalently
its Cartier boundary divisor given by these heights is relatively
ample, with the sign corresponding to the lower-hull convention.
The monomial local charts are the same charts used in Hartl,
Proposition 2.1. This description selects a particular ordered
staircase resolution, not every possible order of Hartl blowups.

\begin{lemma}[Binary component tower, including triple times triple]
\label{r20px:tower}
The natural morphism
\[
                f_{ij}:P_{ij}\longrightarrow S_i\times T_j
 \tag{PX20.2}
\]
is a sequence of blowups of smooth codimension-two centers. The
centers are strict transforms of the products $E_{ia}\times F_{jb}$
with either $a<i,b<j$ or $a>i,b>j$. Lower pairs are ordered in
ascending product order and upper pairs in descending product order.
Every center is a smooth projective surface birational to a product
of two smooth projective curves. Consequently
\[
 h(P_{ij})=h(S_i\times T_j)\oplus K_{ij},
 \qquad K_{ij}\text{ is of abelian type},
 \tag{PX20.3}
\]
with actual rational Chow correspondences for the decomposition and
its inverse. The projector onto the base summand is
$f_{ij}^{*}f_{ij*}$.
\end{lemma}

\begin{proof}
Here is the complete fan calculation. For a positive orthant with
basis $u_1,\ldots,u_m,v_1,\ldots,v_n$, set $u_0=v_0=0$.
Star-subdivide each two-face $\langle u_a,v_b\rangle$ by
$u_a+v_b$, in any reverse-product order. The maximal cones become
\[
 \langle u_a+v_b:(a,b)\in C\setminus\{(0,0)\}\rangle,
 \tag{PX20.4}
\]
where $C$ runs through monotone paths from $(0,0)$ to $(m,n)$.
Proof is by induction on $m+n$. The first ray is $u_m+v_n$.
Its two cones omit respectively $v_n$ and $u_m$. Subsequent
subdivisions restrict to the $(m,n-1)$ and $(m-1,n)$ problems,
each with the extra final ray. A different two-face subdivision
cannot remove the two-face $\langle u_a,v_b\rangle$ before its
own step. Each center is therefore a smooth codimension-two
toroidal stratum. The induction gives exactly the paths classified
by their last step.

The quotient star of $n_{ij}$ in (PX20.1) has coordinates
\[
 n_{aj}\mapsto u_a,\quad n_{ib}\mapsto v_b,\quad
 n_{ab}\mapsto u_a+v_b\quad(a\ne i,b\ne j).
 \tag{PX20.5}
\]
The last formula follows from
$n_{ab}+n_{ij}=n_{aj}+n_{ib}$. A path through $(i,j)$ splits
uniquely into the lower and upper paths; their coordinate sets are
disjoint. Apply (PX20.4) to those two orthants, reversing the lower
orders. This gives precisely the center list and order asserted in
the lemma. For example, the star at $(0,0)$ of
$\Delta^2\times\Delta^2$ is obtained by the four subdivisions
\[
 u_2+v_2,\quad u_2+v_1,\quad u_1+v_2,\quad u_1+v_1.
 \tag{PX20.6}
\]
It has the six monotone-path maximal cones. No triple stratum is
removed or replaced by a node assumption.

Globalize this list by the fixed global orders. Restricting it to
indices present in a chart gives precisely the local list just
proved. The centers are intersections of the relevant two strict
transform boundary divisors; their two-face has not yet been
subdivided. The reduced total boundary remains SNC after every
step, proving smoothness of the centers. Both modifications have
the same local fan and are the identity on their common dense open;
the identifications therefore glue over $S_i\times T_j$.

Each center has a projective birational map to its original
$E_{ia}\times F_{jb}$. A birational morphism of smooth projective
surfaces factors into point blowups, in arbitrary characteristic
(Stacks Project, Lemma 54.17.1, Tag 0C5R).
The curve-product motive and the additional point motives are of
abelian type. The smooth blowup motive formula proves (PX20.3)
with the indicated base projector. This argument does not assume
that closed strata of $\mathcal P$ are toric bundles.
\end{proof}

\subsection{Actual same-label double strata}

\begin{lemma}[No hidden double stratum]
\label{r20px:double}
Every connected double stratum joining $P_{ij}$ to $P_{il}$ maps
properly birationally to a connected component of
$S_i\times (T_j\cap T_l)$ and is an isomorphism away from
$\partial S_i\times (T_j\cap T_l)$. There is no additional
double component supported entirely over $\partial S_i$.
If the $S$ label changes instead, the double stratum maps into
$\partial S_i$ on its $S_i$ side. The symmetric assertions hold
with $S,T$ interchanged.
\end{lemma}

\begin{proof}
Over $S_i^{\circ}=S_i\setminus\partial S_i$, the first local
simplex is $\Delta^0$. The staircase modification is the ordinary
product, giving the asserted dense open map. To exclude an extra
component, let $x_a$ be the local equation of a different first-factor
component. Its valuation is zero on both rays $n_{ij},n_{il}$.
It is therefore invertible at the generic orbit of their two-cone.
Thus a double component cannot be supported over $x_a=0$.
Equivalently, its support there would force it into a third distinct
vertical component, contradicting the SNC codimension of a triple
intersection. The same argument holds on all local charts and after
passing to each connected component. Each remaining-factor double
component therefore has exactly one dense open inverse image.
Properness gives the birational statement. When the first labels
differ, the map factors through their intersection by definition.
\end{proof}

\subsection{A supported-cohomology transport calculation}

\begin{lemma}[Transport with all exceptional inputs]
\label{r20px:transport}
Let $f:D\to B$ be a proper birational morphism of smooth projective
varieties, an isomorphism off $Z\subset B$. Let $p$ be a self-adjoint
algebraic cohomological projector on $B$ annihilating
$\operatorname{im}(H_Z^*(B)\to H^*(B))$. Then $P=f^*pf_*$ is
a self-adjoint algebraic projector. In a commutative square
\[
\begin{matrix}
 E&\xrightarrow{g}&D\\
 f'\downarrow&&\downarrow f\\
 B'&\xrightarrow{h}&B,
\end{matrix}
 \tag{PX20.7}
\]
assume $g,h$ are smooth-divisor embeddings, $f'$ is proper
birational and an isomorphism off $Z'=h^{-1}Z$, and $p'$ is the
analogous projector on $B'$ satisfying $h^*p=p'h^*$. Then
\[
 g^*P=P'g^*,\qquad Pg_*=g_*P',\qquad P'=f'^*p'f'_*.
 \tag{PX20.8}
\]
\end{lemma}

\begin{proof}
The identity $f_*f^*=1$ gives idempotence, and transposition gives
self-adjointness. Decompose $\beta\in H^*(D)$ as
$f^*f_*\beta+\beta_0$, with $f_*\beta_0=0$. On the first
summand (PX20.8) follows from the square and $h^*p=p'h^*$.
The class $\beta_0$ vanishes on the locus where $f$ is an
isomorphism. Localization puts it in supported cohomology over
$f^{-1}Z$. Its pullback to $E$ and pushforward to $B'$ are
supported over $Z'$. The projector $p'$ kills the latter class.
This proves the first identity on the exceptional summand as well.
Taking Poincare adjoints gives the second identity.
\end{proof}

For a smooth projective surface $S$, choose the self-adjoint
surface Chow--Kunneth projector $p_S$ onto the orthogonal complement
of divisor classes in $H^2(S)$, zero in other degrees. This is an
actual rational algebraic projector; see Kahn--Murre--Pedrini,
Propositions 7.2.1 and 7.2.3. For a union $B$ of SNC boundary curves,
\[
 p_S\operatorname{im}(H_B^*(S)\to H^*(S))=0,
 \qquad p_S|_{H^*(B)}=0
 \tag{PX20.9}
\]
where the second formula means restriction after $p_S$.
Indeed the only supported image in degree two is generated by the
boundary divisor classes; point intersections contribute nothing
in that degree. Restriction to a boundary curve is the same divisor
pairing. Kunneth with supports gives the same assertions for
$p_S\boxtimes1$ and $B\times T$, for any smooth projective $T$.
Only these cohomological annihilations are used; no Chow-level
annihilation by arbitrary maps to curves is inferred.

\subsection{The four sectors and a cohomological inverse}

Let $A=\sigma\rho$ be the actual incidence correspondence of
$\mathcal P$, of degree one. On each $P_{ij}$ put
\[
 P^S=f_{ij}^*(p_{S_i}\boxtimes1)f_{ij*},\qquad
 P^T=f_{ij}^*(1\boxtimes p_{T_j})f_{ij*}.
 \tag{PX20.10}
\]
They commute because $f_{ij*}f_{ij}^*=1$. Keep individual
$S_i$ and $T_j$ labels separate on the direct sum of components.
Lemma \ref{r20px:double}, (PX20.9), and
Lemma \ref{r20px:transport} show that every boundary restriction
and its Gysin adjoint commute with the corresponding projectors.
An edge changing a transcendental factor label has zero restriction
on that sector. Thus $A$ preserves the four orthogonal sectors
\[
 P^SP^T,\quad P^S(1-P^T),\quad(1-P^S)P^T,\quad
 Q=(1-P^S)(1-P^T).
 \tag{PX20.11}
\]

On $P^SP^T$, $A=0$. On the summand with fixed $t_2(S_i)$,
transport back by $f_{ij*},f_{ij}^*$ identifies $A$ with
\[
                    1_{t_2(S_i)}\boxtimes A_{\mathcal T}.
 \tag{PX20.12}
\]
For an off-diagonal entry this follows from the degree-one
birational map of Lemma \ref{r20px:double}: its pushforward after
pullback is the identity. The diagonal is the negative sum of the
boundary self-intersections; applying the same argument with the
same endpoint on both sides includes every diagonal term. The
entries changing the first label are zero. The assertion with
$S,T$ reversed follows identically.

Recall why the required surface inverse is already rigorous without
surface Tate conjectures. On the disjoint union of native surface
components the incidence is a correspondence in $\operatorname{CH}^3$.
Jannsen's numerical semisimplicity gives a numerical inner inverse
represented by a divisor correspondence $G\in\operatorname{CH}^1$.
The error $A-AGA$ is a numerical-zero curve on a product of two
surfaces. Frozen SP19 curve comparison makes its entire geometric
class zero. Hence $AGA=A$ on full cohomology for every invertible
prime. To use it on $1-p_T$ in (PX20.11), replace $G$ by
$(1-p_T)G(1-p_T)$; $A$ annihilates $p_T$ on both sides, so the
inner-inverse identity persists. This splits both single-$t_2$
sectors of (PX20.11).

By Lemma \ref{r20px:tower}, the remaining $Q$ sector is of abelian
type: its base is the tensor product of the nontranscendental surface
parts, and its additional summand is $K_{ij}$. The former parts
consist of Tate and Picard/Albanese motives. Choose one finite product
of abelian varieties containing all these finitely many motives and
all needed correspondence groups inside suitable finite sums and
Tate shifts of its powers. Theorem 1 and Remark (b) of
Deligne--Clozel (2009) supply a single invertible prime $\ell$ for
which numerical and geometric homological equivalence coincide on
all those powers. Jannsen's numerical semisimplicity then provides
an algebraic inner inverse for $QAQ$, including each degree and
Tate block. This does not apply numerical semisimplicity to an
unidentified full category of homological motives.
For the native residue-field version, numerical zero remains zero
against geometric test cycles: descend any such test to a finite
constant extension and use its norm, the full extension degree,
and the projection formula. Thus native numerical-zero
correspondences also have zero geometric class at this selected
prime. No degree on the absolute arithmetic model is used.

Summing the three nonzero-sector inverses gives an actual rational
correspondence $G$ on $\coprod P_{ij}$ with
\[
                         A^HG^HA^H=A^H
 \tag{PX20.13}
\]
on all geometric cohomology, at the one selected $\ell$.
All choices over a finite residue extension descend by normalized
Galois averaging: (PX20.13) is linear in $G$, and $A$ descends.
The claim here is at one suitable prime; an all-prime assertion for
the residual abelian-type core has not been proved by this argument.

\begin{theorem}[Two bad surfaces with arbitrary good coefficients]
\label{r20px:local}
Assume the frozen AK19, OW19, VT17, FD17 and R16 descent interfaces.
Fix a finite valued extension $K/\mathbf Q_p$ and its original
absolute value. Let $S,T/K$ be arbitrary smooth projective
geometrically integral surfaces. Let $V/K$ be any smooth projective
geometrically integral variety having a smooth projective model
after a finite valued extension. Then every specified geometrically
homologically trivial rational cycle on $V\times S\times T$, in
every codimension $1\leq r\leq\dim(V\times S\times T)$, has an
ordinary-actual strong Green form on the
complement of its exact support, tested against all actual compact
forms. The assertion applies afresh after every finite valued
extension and to all new cycles over that extension.
\end{theorem}

\begin{proof}
Take a common finite valued extension giving a smooth projective
model of $V$ and semistable projective alterations of both surfaces.
Repair further ramified factor base changes by Hartl 2.2 before
forming the product, as in frozen SP19. Split component constants
unramifiedly and use the ordered staircase model above.
Its product with the smooth model of $V$ has incidence and inverse
$\Delta_{V_0}\boxtimes A$ and $\Delta_{V_0}\boxtimes G$.
Equation (PX20.13) holds on full cohomology, including all mixed
Kunneth summands. Both maps are actual Chow correspondences and
therefore preserve the actual algebraic class span and its
numerical projection. The correction $x-G^HA^Hx$ proves
$\Delta_A=0$ in every degree, exactly as in PI20/AK19.

All degrees of surface cohomology satisfy weight--monodromy, and
the good factor has pure unramified cohomology. Kunneth and tensor
stability give the odd-weight hypothesis of frozen OW19 in every
required degree. Thus the actual inertia image is zero for this
specific witness. AK19 makes every pulled-back homological cycle
obstruction zero. R16 dominant detection and finite trace descent,
with the full field degree and the original absolute value, return
this vanishing to the specified original cycle; VT17 gives the
asserted actual form with all tests and exact support. No alteration
is assumed finite on the specified complement.
\end{proof}

\paragraph{Exact boundary of the present verification.}
The statement for two surfaces alone is already contained in frozen
SP19's four-dimensional all-codimension sector. The extra good
coefficient allows higher-dimensional mixed cycles, including the
middle codimension-three sector when $\dim V=1$. The idea and its
historical candidate proof occurred in R9 ST/TR; this note makes no
priority claim and does not silently revise historical registers.
For arbitrary products of three or more bad surfaces, the deletion
projector and residual-core argument requires additional geometric
generator compatibility. The R9 DG/GL/DT files propose such a proof,
but this binary verification alone does not establish it. General
strong $G$ and arbitrary-variety complete Beilinson--Bloch height
remain open.

\paragraph{Primary sources checked for scope.}
Hartl, \emph{Semi-stability and base change}, Arch. Math. 77 (2001),
215--221, Proposition 2.1 and its two monomial charts on 218--219,
and Proposition 2.2 on 219--221. Kahn--Murre--Pedrini,
\emph{On the transcendental part of the motive of a surface},
July 2005 author manuscript, Propositions 7.2.1 and 7.2.3,
pp.12--15; the source uses covariant motives, so only its stated
projectors and realized orthogonality are imported here.
Deligne (written by Clozel), \emph{Equivalence numerique,
equivalence cohomologique, et theorie de Lefschetz des varietes
abeliennes sur les corps finis}, PAMQ 5 (2009), 495--506,
Theorem 1 and Remarks (a)--(b), p.496. Jannsen,
\emph{Motives, numerical equivalence, and semi-simplicity},
Invent. Math. 107 (1992), Theorem 1, used only in numerical motives.
The surface factorization input is Stacks Project, Tag 0C5R;
the blowup and projective-bundle Chow formulas are Tags 0FVA and
0FGQ, applied after every smooth projective test product to obtain
the Chow-motive identity (Tag 0FGK, Manin's identity principle).
No assertion in these references by itself supplies the component
tower or projector transport: those are proved above.
